\documentclass[../../../include/open-logic-section]{subfiles}

\begin{document}

\olfileid{mod}{mar}{mpa}
\section{$\Th{PA}$的模型}

% Part: first-order-logic
% Chapter: models-of-arithmetic
% Section: models-of-pa

\begin{explain}
任意一个~$\Th{TA}$ 的非标准模型也是~$\Th{PA}$ 的模型。
我们知道~$\Th{TA}$（从而~$\Th{PA}$）的非标准模型是存在的。
我们还知道，这样的非标准模型包含非标准“数”，
即论域中“超越”所有标准“数”的元素。
但是，它们是如何排列的？有多少个呢？
我们看到过，更弱理论~$\Th{Q}$ 的模型可以只包含一个非标准数。
但这些简单的结构并不是 $\Th{PA}$ 或 $\Th{TA}$ 的模型。

理解 $\Th{PA}$ 或 $\Th{TA}$ 模型结构的关键，
在于考察这些理论能推导出哪些事实。
例如，光是 $\Th{PA}$ 就能证明 $\lforall[x][\eq/[x][x']]$
和 $\lforall[x][\lforall[y][\eq[(x+y)][(y+x)]]]$，这就排除了（使这些语句为假的）简单结构作为~$\Th{PA}$ 模型的可能性。

假设~$\Struct{M}$ 是~$\Th{PA}$ 的一个模型。那么，若 $\Th{PA} \Proves !A$，就有 $\Sat{M}{!A}$。
我们再次用 $\nszero$ 表示 $\Assign{\Obj{0}}{M}$，
用 $\nssucc$ 表示 $\Assign{\prime}{M}$，
用 $\nsplus$ 表示 $\Assign{+}{M}$，
用 $\nstimes$ 表示 $\Assign{\times}{M}$，
并用 $\nsless$ 表示 $\Assign{<}{M}$。
于是，任意语句~$!A$ 都陈述了关于 $\nszero$、$\nssucc$、$\nsplus$、$\nstimes$ 和 $\nsless$ 的某个条件，而如果 $\Sat{M}{!A}$，则该条件必须被满足。
例如，如果 $\Sat{M}{!Q_1}$，即
$\Sat{M}{\lforall[x][\lforall[y][(\eq[x'][y'] \lif \eq[x][y])]]}$，
那么 $\nssucc$ 必须是单射的。
\end{explain}

\begin{prop}
在 $\Struct{M}$ 中，$\nsless$ 是一个严格线序，即满足：
\begin{enumerate}
\item 对所有 $x \in \Domain{M}$，$x \nsless x$ 不成立。
\item 若 $x \nsless y$ 且 $y \nsless z$，则 $x \nsless z$。
\item 对任意 $x \neq y$，或者 $x \nsless y$，或者 $y \nsless x$。
\end{enumerate}
\end{prop}

\begin{proof}
$\Th{PA}$ 证明了：
\begin{enumerate}
\item $\lforall[x][\lnot x < x]$
\item $\lforall[x][\lforall[y][\lforall[z][((x < y \land y < z) \lif x < z)]]]$
\item $\lforall[x][\lforall[y][((x < y \lor y < x) \lor \eq[x][y]))]]$
\end{enumerate}
\end{proof}

\begin{prop}
\ollabel{prop:M-discrete}
在 $\nsless$-序下，$\nszero$ 是 $\Domain{M}$ 的最小元。
对任意 $x$，有 $x \nsless x^\nssucc$，并且 $x^\nssucc$ 是具有该性质的 $\nsless$-最小元。
对任意 $x$，存在唯一的 $y$ 使得 $y^\nssucc = x$。（我们称 $y$ 为 $x$ 在 $\Struct{M}$ 中的“前驱”，并记作 $^\nssucc x$。）
\end{prop}

\begin{proof}
练习。
\end{proof}

\begin{prob}
请在~$\Lang{L_A}$ 中找出一些语句，使得它们在~$\Th{PA}$ 中可推导（从而在~$\Struct{N}$ 中为真），并且能保证
\olref[mod][mar][mpa]{prop:M-discrete} 中 $\nszero$、$\nssucc$ 和 $\nsless$ 的性质。
\end{prob}

\begin{prop}
$\Struct{M}$ 的所有标准元素都小于（按照~$\nsless$）所有非标准元素。
\end{prop}

\begin{proof}
我们用 $n$ 简记 $\Value{\num{n}}{M}$，即 $\Struct{M}$ 的一个标准元素。$\Th{Q}$ 已经证明，对于任何 $n \in \Nat$，$\lforall[x][(x < \num{n}' \lif (\eq[x][\num{0}] \lor
  \eq[x][\num{1}] \lor \dots \lor \eq[x][\num{n}]))]$。不存在小于 $\nszero$ 的元素。所以，如果 $n$ 是标准的而 $x$ 是非标准的，我们不可能有 $x \nsless n$。根据定义，非标准元素是指对任何 $n \in \Nat$ 都不等于 $\Value{\num{n}}{M}$ 的元素，因此也有 $x \neq n$。由于 $\nsless$ 是一个线序，我们必然有 $n \nsless x$。
\end{proof}

\begin{prop}
$\Domain{M}$ 的每个非标准元素~$x$ 都属于子集 \[
\dots ^{\nssucc\nssucc\nssucc}x \nsless ^{\nssucc\nssucc}x \nsless
^{\nssucc}x \nsless x \nsless x^{\nssucc} \nsless x^{\nssucc\nssucc}
\nsless x^{\nssucc\nssucc\nssucc} \nsless \dots
\]
我们称这个子集为 $x$ 的\emph{块}，记作 $[x]$。它既没有最小元也没有最大元。它可以刻画为这样的 $y \in \Domain{M}$ 的集合：存在某个标准~$n$，使得 $x \nsplus n = y$ 或 $y \nsplus n = x$。
\end{prop}

\begin{proof}
显然，这样的集合 $[x]$ 总是存在的，因为 $\Domain{M}$ 的每个元素 $y$ 都有唯一后继元 $y^\nssucc$ 和唯一前驱元 $^\nssucc y$。对于相继的元素 $y$、$y^\nssucc$，我们有 $y \nsless y^\nssucc$，并且 $y^\nssucc$ 是 $\Domain{M}$ 中满足 $y$ 小于它的 $\nsless$-最小元。由于总是 $^\nssucc y \nsless y$ 且 $y \nsless y^\nssucc$，$[x]$ 既没有最小元也没有最大元。如果 $y \in [x]$，那么 $x \in [y]$，因为此时要么 $y^{\nssucc\dots\nssucc} = x$，要么 $x^{\nssucc\dots\nssucc} = y$。如果 $y^{\nssucc\dots\nssucc} = x$（使用了 $n$ 个 $\nssucc$），那么 $y \nsplus n = x$，反之亦然，因为 $\Th{PA} \Proves \lforall[x][\eq[x^{\prime\dots\prime}][(x +
    \num{n})]]$（若 $n$ 是 $\prime$ 的个数）。
\end{proof}

\begin{prop}
如果 $[x] \neq [y]$ 且 $x \nsless y$，那么对于任意 $u \in [x]$ 和任意 $v \in [y]$，都有 $u \nsless v$。
\end{prop}

\begin{proof}
注意$\Th{PA} \Proves \lforall[x][\lforall[y][(x < y \lif (x' < y
    \lor x' = y))]]$。因此，如果 $u \nsless v$，并且 $[u] \neq [v]$，那么对任意~$n$ 也有 $u \nsplus n^\nssucc \nsless v$。

任意 $u \in [x]$ 都满足 $u \nsless y$：根据假设有 $x \nsless y$。如果 $u \nsless x$，由传递性得 $u \nsless y$。而如果 $x \nsless u$ 但 $u \in [x]$，则对某个~$n$ 有 $u = x\nsplus n^\nssucc$，从而由刚才证明的事实得 $u \nsless y$。

现在假设 $v \in [y]$ 且 $v \nsless y$，即对某个标准~$m$ 有 $v \nsplus m^\nssucc = y$。这排除了 $v \nsless x$ 的可能性，否则 $y = v \nsplus m^\nssucc \nsless x$。显然也有 $x \neq v$，否则$x \nsplus m^\nssucc = v \nsplus m^\nssucc = y$，我们将得到 $[x] = [y]$。所以，$x \nsless v$。但这样一来，对任意~$n$ 也有 $x \nsplus n^\nssucc \nsless v$。因此，如果 $x \nsless u$ 且 $u \in [x]$，则有 $u \nsless v$。如果 $u \nsless x$，那么由传递性得 $u \nsless v$。

最后，如果 $y \nsless v$，由于我们已经证明 $u \nsless y$ 且 $y \nsless v$，所以 $u \nsless v$。
\end{proof}

\begin{cor}
如果 $[x] \neq [y]$，则 $[x] \cap [y] = \emptyset$。
\end{cor}

\begin{proof}
假设 $z \in [x]$ 且 $x \nsless y$。那么对所有 $u \in [y]$ 有 $z \nsless u$。如果 $z \in [y]$，我们将得到 $z \nsless z$。若 $y \nsless x$，同理可得。
\end{proof}

\begin{explain}
这意味着块本身可以以一种保持 $\nsless$ 的方式排序：$[x] \nsless [y]$ 当且仅当 $x \nsless y$，或者等价地说，对任意 $u \in [x]$ 和 $v \in [y]$ 都有 $u \nsless v$。显然，标准块 $[0]$ 是最小的块。它与任何非标准块都不相交，并且任意两个非标准块也不相交。具体来说，你无法通过重复取后继或前驱来“到达”一个不同的块。
\end{explain}

\begin{prop}
如果 $x$ 和 $y$ 都是非标准的，那么 $x \nsless x \nsplus y$ 且 $x \nsplus y \notin [x]$。
\end{prop}

\begin{proof}
如果 $y$ 是非标准的，那么 $y \neq \nszero$。$\Th{PA} \Proves
\lforall[x][(y \neq \Obj{0} \lif x < (x+y))]$。现在假设 $x \nsplus y \in [x]$。由于 $x \nsless x \nsplus y$，我们将有 $x \nsplus n^\nssucc = x \nsplus y$。但是$\Th{PA} \Proves
\lforall[x][\lforall[y][\lforall[z][(\eq[(x+y)][(x+z)] \lif y = z)]]]$（加法的消去律）。这意味着对某个标准~$n$ 有 $y = n^\nssucc$；但 $y$ 被假定为非标准的。
\end{proof}

\begin{prop}
不存在最小的非标准块。
\end{prop}

\begin{proof}
$\Th{PA} \Proves \lforall[x][\lexists[y][(\eq[(y+y)][x] \lor
      \eq[(y+y)'][x])]]$，即每个 $x$ 都可被~$2$ 整除（可能余~$1$）。如果 $x$ 是非标准的，那么~$y$ 也是非标准的。根据前面的命题，$y \nsless y \nsplus y$ 且 $y \nsplus y \notin [y]$。那么也有 $y \nsless (y \nsplus y)^\nssucc$ 且$(y
  \nsplus y)^\nssucc \notin [y]$。但是 $x = y \nsplus y$ 或 $x = (y \nsplus y)^\nssucc$，所以 $y \nsless x$ 且 $y \notin [x]$。
\end{proof}

\begin{prop}
不存在最大的块。
\end{prop}

\begin{proof}
练习。
\end{proof}

\begin{prob}
证明在~$\Th{PA}$ 的非标准模型中，不存在最大的块。
\end{prob}

\begin{prop}
\ollabel{prop:blocks-dense}
块的序是稠密的。也就是说，如果 $x \nsless y$ 且 $[x] \neq [y]$，那么存在一个不同于两者的块 $[z]$ 介于它们之间。
\end{prop}

\begin{proof}
假设 $x \nsless y$。如前所述，$x \nsplus y$ 可被 2 整除（可能有余数）：存在 $z \in \Domain{M}$ 使得要么 $x \oplus y = z \oplus z$，要么$x \oplus y = (z \oplus
z)^\nssucc$。元素 $z$ 是 $x$ 和~$y$ 的“平均”，并且 $x \nsless z$ 且 $z \nsless y$。
\end{proof}

\begin{prob}
详细写出 \olref[mod][mar][mpa]{prop:blocks-dense} 的证明。$\Th{PA}$ 必须推导出哪个语句才能保证~$z$ 的存在？为什么 $x \nsless z$ 且 $z \nsless y$，以及为什么 $[x] \neq [z]$ 且 $[z] \neq [y]$？
\end{prob}

\begin{explain}
因此，非标准块的排序方式就像有理数一样：它们构成一个可数无穷、稠密、线性且无端点的序。可以证明，任意两个这样的可数序都是同构的。由此可得，对于算术真理论的任意两个可数非标准模型 $\Struct{M}_1$ 和 $\Struct{M_2}$，它们到只包含 $<$ 和 $=$ 的语言的归约是同构的。事实上，我们可以如下定义一个同构 $h$：$\Struct{M}_1$ 和 $\Struct{M}_2$ 的标准部分都与标准模型 $\Struct{N}$ 同构，因此彼此也同构。构成非标准部分的那些块本身也像有理数一样排序，从而是同构的；我们可以将块之间的同构扩展为\emph{块内}的同构，方法是在每个块中任意匹配一个元素，然后规定 $\Struct{M}_1$ 中 $x$ 的后继在 $\Struct{M}_2$ 中的像，就是 $x$ 的像的后继。需注意，这\emph{并不}意味着 $\mathfrak{M}_1$ 和 $\mathfrak{M}_2$ 在全算术语言下是同构的（事实上，同构总是相对于某个语言而言的），因为在 $\Domain{M_1}$ 和 $\Domain{M_2}$ 上定义加法和乘法可以有不同的、不同构的方式。（这也源自 Vaught 的一个著名定理：一个完备理论的可数模型个数不可能是 2。）
\end{explain}

\end{document}